\documentclass[10pt,a4paper]{amsart}
\usepackage[margin=19mm]{geometry}
\usepackage{amsmath,amssymb,mathtools}
\usepackage[scr=rsfs,cal=euler]{mathalfa}
\usepackage{xfrac}
\usepackage[unicode,hidelinks,bookmarks=false]{hyperref}
\newcommand{\ie}{i.e.\ }
\newcommand{\cf}{cf.\ }
\newcommand{\resp}{respectively\ }
\newcommand{\Subsection}{\subsection}
\providecommand{\nobreakdash}{\nobreak\hbox{-}}
\setlength{\emergencystretch}{2em}
\multlinegap0pt
\usepackage{amssymb}
\usepackage[shortlabels]{enumitem}
\newtheorem{proposition}{Proposition}
\newcommand\Lc{{\mathcal{L}}}
\newcommand\Tr{\mathrm{Tr}}
\newcommand\U{{\mathrm U}}
\title{Universal dualities for Wilson loops in lattice Yang--Mills}
\author{Thibaut Lemoine}
\date{Source: arXiv:2604.16252v2 (CC BY 4.0)}
\begin{document}
\maketitle

\noindent\emph{Source and licence.} Universal dualities for Wilson loops in lattice Yang--Mills. Version arXiv:2604.16252v2 (CC BY 4.0), distributed by arXiv under the Creative Commons Attribution 4.0 International licence, http://creativecommons.org/licenses/by/4.0/, as recorded in the arXiv licence metadata for this exact version and shown on the abstract page https://arxiv.org/abs/2604.16252v2. Full licence notice: \texttt{licenses/arxiv-2604.16252-CC-BY-4.0.txt}.\par\smallskip

\noindent\textit{Editorial scope.} Counted unit: the article's proposition prop:universal-loop-laplacian-term (source0.tex lines 1865--1871). The article's complete original proof of it is delivered with it (source0.tex lines 1873--1990, one \texttt{proof} environment of 118 lines). No mathematics is written, repaired or re-derived: the statement and the proof are the article's own lines, in the article's own notation, and no retained line is substituted or re-flowed.  Retained uncounted as the source-local context the proof needs: lines 1637--1639.  Nothing was omitted from any retained span, so the package carries no omission record.  No retained line contains a citation, so the wrapper carries no bibliography and no bibliography entry.  No retained line contains a figure environment, an image-inclusion command or drawing code, so no image is delivered and the package declares no asset dependency.  Editorial formatting only, in addition: an AMS \texttt{amsart} preamble that restates the article's own declarations and macros used by the retained lines, namely the environments \texttt{proposition} on separate counters, together with the standard packages those lines need.  The retained environments print in this document as Proposition 1 in order of appearance, whereas the article prints the same items with the numbering of its own full document.  The complete native source of the article as distributed in the arXiv e-print for this exact version is bundled unchanged as \texttt{sources/198672/source0.tex}.
\begin{equation}\label{eq:magic}
\sum_{a=1}^{N^2}X_a^2=-NI_N,\qquad \sum_{a=1}^{N^2} X_aAX_a=-\Tr(A)I_N,\qquad \sum_{a=1}^{N^2} \Tr(X_aA)\Tr(X_aB)=-\Tr(AB).
\end{equation}

\begin{proposition}
\label{prop:universal-loop-laplacian-term}
For every oriented edge $e$, loop sequence $\Lc$, and plaquette decoration $\alpha$, one has
\[
I_\Delta:=\int_{\U(N)^{E(\Lambda)}}(\Delta_e W_{\Lc}(U))\prod_{p\in P(\Lambda)}\chi_{\alpha_p}(U_{\partial p}) dU=(\mathscr L_e \widehat{W}_{\Lambda,\Lc})(\alpha).
\]
\end{proposition}

\begin{proof}
We prove first a pointwise identity for $\Delta_eW_{\Lc}$, and then integrate against the undifferentiated plaquette factors $\prod_{p\in P(\Lambda)}\chi_{\alpha_p}(U_{\partial p})$. Fix an oriented edge $e\in E(\Lambda)$. For each $i\in\{1,\dots,k\}$, choose the cyclic representative
\[
w_{\ell_i}=e_{i,1}^{\varepsilon_{i,1}}\cdots e_{i,m_i}^{\varepsilon_{i,m_i}},\qquad \varepsilon_{i,r}\in\{\pm1\}.
\]
If $x\in A_i(e;\Lc)$, write $w_{\ell_i}=a_xeb_x,$ and if $x\in B_i(e;\Lc)$, write $w_{\ell_i}=a_xe^{-1}b_x.$ By differentiation, we have
\[
\partial_{e,a}\Tr(U_{\ell_i})=\sum_{x\in A_i(e;\Lc)}\Tr\bigl(U_{a_x}X_aU_eU_{b_x}\bigr)-\sum_{x\in B_i(e;\Lc)}\Tr\bigl(U_{a_x}U_e^{-1}X_aU_{b_x}\bigr).
\]
Using cyclicity of the trace, this may be rewritten as
\begin{equation}\label{eq:local-first-derivative-loop}
\partial_{e,a}\Tr(U_{\ell_i})=\sum_{x\in A_i(e;\Lc)} \Tr\bigl(X_aU_{e b_x a_x}\bigr)-\sum_{x\in B_i(e;\Lc)} \Tr\bigl(X_aU_{b_x a_x e^{-1}}\bigr).
\end{equation}

Since $W_{\Lc}=\prod_{i=1}^k \Tr(U_{\ell_i})$, the first derivative is
\[
\partial_{e,a}W_{\Lc}=\sum_{i:e\in D_i(\Lc)}\left(\prod_{j\neq i}\Tr(U_{\ell_j})\right)\partial_{e,a}\Tr(U_{\ell_i}),
\]
and therefore, by differentiating once more and summing over $a$,
\[
\Delta_eW_{\Lc}=\sum_{i=1}^k\left(\prod_{j\neq i}\Tr(U_{\ell_j})\right)\Delta_e\Tr(U_{\ell_i})+\sum_{i\neq j}\sum_{a=1}^{N^2}\left(\prod_{m\neq i,j}\Tr(U_{\ell_m})\right)\bigl(\partial_{e,a}\Tr(U_{\ell_i})\bigr)\bigl(\partial_{e,a}\Tr(U_{\ell_j})\bigr).
\]
We analyze separately the two terms on the right-hand side and use the magic formulas~\eqref{eq:magic}.

\smallskip

\emph{Step 1: second derivatives inside a single loop.}
Fix $i$ with $e\in D_i(\Lc)$. Differentiating \eqref{eq:local-first-derivative-loop} once more, the two derivatives may hit either the same active occurrence or two distinct active occurrences.
\begin{itemize}
\item If they hit the same active occurrence $x\in C_i(e;\Lc)$, then one gets a factor $\sum_a X_a^2$, hence, with the normalization matching the definition of $\Lc_e$,
\[
\sum_{a=1}^{N^2}\partial_{e,a}^2 \Tr(U_{\ell_i})\Big|_{\text{same }x}=-N\Tr(U_{\ell_i}).
\]
Summing over the $m_i(e;\Lc)=|C_i(e;\Lc)|$ active occurrences gives the contribution
\[
-Nm_i(e;\Lc)\Tr(U_{\ell_i}).
\]
\item If they hit two distinct active occurrences $x\neq y\in C_i(e;\Lc)$, we have to distinguish between four cases.
\begin{enumerate}
\item If $x,y\in A_i(e;\Lc)$, write $w_{\ell_i}=aebec$. Then the contribution is
\[
\sum_{a=1}^{N^2}\Tr(U_aX_aU_eU_bX_aU_eU_c).
\]
By cyclicity of the trace and the standard $\U(N)$ contraction identity in the chosen
normalization,
\[
\sum_{a=1}^{N^2}\Tr(AX_aBX_a)=-\Tr(A)\Tr(B),
\]
this equals $-\Tr(U_{aec})\Tr(U_{be}),$ which is exactly the Wilson weight associated with the positive splitting $(\times^1_{x,y}\ell_i,\times^2_{x,y}\ell_i)$.
\item If $x,y\in B_i(e;\Lc)$, writing $w_{\ell_i}=ae^{-1}be^{-1}c$, the same computation gives $-\Tr(U_{aec})\Tr(U_{be^{-1}}),$
 i.e. the positive splitting contribution again.
\item If $x\in A_i(e;\Lc)$ and $y\in B_i(e;\Lc)$, write $w_{\ell_i}=aebe^{-1}c$. Then the second
derivative carries one minus sign from differentiating the inverse occurrence, and one gets
\[
-\sum_{a=1}^{N^2}\Tr(U_aX_aU_eU_bU_e^{-1}X_aU_c).
\]
Using cyclicity and the same contraction identity, this becomes $\Tr(U_{ac})\Tr(U_b),$ which is exactly the contribution of the negative splitting $(\times^1_{x,y}\ell_i,\times^2_{x,y}\ell_i)$.
\item If $x\in B_i(e;\Lc)$ and $y\in A_i(e;\Lc)$, the case is identical to (iii) and gives the other half of
$S^-_{i,e}(\Lc)$.
\end{enumerate}
\end{itemize}

Collecting all same-loop contributions, we obtain
\begin{align}
\Delta_e\Tr(U_{\ell_i}) &=-Nm_i(e;\Lc)\Tr(U_{\ell_i})+ \sum_{\Lc'\in S^-_{i,e}(\Lc)} W_{\Lc'}^{(i)}- \sum_{\Lc'\in S^+_{i,e}(\Lc)} W_{\Lc'}^{(i)},\label{eq:single-loop-Laplacian}
\end{align}
where $W_{\Lc'}^{(i)}$ means the product of traces obtained from $\Lc'$ after reinserting the undifferentiated loops $\ell_j$, $j\neq i$.

\smallskip

\emph{Step 2: cross terms between two different loops.}
Fix $i\neq j$ with $e\in D_i(\Lc)\cap D_j(\Lc)$. By \eqref{eq:local-first-derivative-loop}, the cross term is a sum over pairs of active occurrences $x\in C_i(e;\Lc)$, $y\in C_j(e;\Lc)$.

If $x,y$ have the same orientation, say $x\in A_i(e;\Lc)$, $y\in A_j(e;\Lc)$, write $w_{\ell_i}=aeb$ and $w_{\ell_j}=ced$. Then
\[
\sum_{a=1}^{N^2}\Tr(U_aX_aU_eU_b)\Tr(U_cX_aU_eU_d).
\]
Using the contraction identity
\[
\sum_{a=1}^{N^2}\Tr(X_aA)\Tr(X_aB)=-\Tr(AB)
\]
in the normalization matching $\Lc_e$, this becomes
\[
-\Tr(U_{aedceb}),
\]
namely the positive $\U(N)$-merger $\ell_i\oplus_{x,y}\ell_j$.

If $x,y\in B_i(e;\Lc)\times B_j(e;\Lc)$, one obtains similarly the other half of $M^+_{U,i,e}(\Lc)$.

If $x,y$ have opposite orientations, say $x\in A_i(e;\Lc)$, $y\in B_j(e;\Lc)$, there is one minus sign coming from differentiating the inverse occurrence in $\ell_j$, and the same contraction identity yields
\[
\Tr(U_{adcb}),
\]
which is the negative $\U(N)$-merger $\ell_i\ominus_{x,y}\ell_j$. The case $x\in B_i(e;\Lc)$, $y\in A_j(e;\Lc)$ is identical.

Therefore the total cross-term contribution is
\begin{align}
\sum_{a=1}^{N^2}
\bigl(\partial_{e,a}\Tr(U_{\ell_i})\bigr)\bigl(\partial_{e,a}\Tr(U_{\ell_j})\bigr)=\sum_{\Lc'\in M^-_{U,i,e}(\Lc;j)} W_{\Lc'}-\sum_{\Lc'\in M^+_{U,i,e}(\Lc;j)} W_{\Lc'},
\label{eq:cross-loop-Laplacian}
\end{align}
where $M^\pm_{U,i,e}(\Lc;j)$ denotes the part of the merger multiset involving the pair
$(i,j)$.

\smallskip

\emph{Step 3: conclusion.}
Substituting \eqref{eq:single-loop-Laplacian} and \eqref{eq:cross-loop-Laplacian} into the product-rule expansion of $\Delta_eW_{\Lc}$, and summing over all active loops and all active pairs of loops, we obtain the pointwise identity
\[
\Delta_eW_{\Lc}(U)=-N\sum_{i:e\in D_i(\Lc)} m_i(e;\Lc)W_{\Lc}(U)+\sum_{i:e\in D_i(\Lc)}\sum_{\Lc'\in S^-_{i,e}(\Lc)} W_{\Lc'}(U)-\sum_{i:e\in D_i(\Lc)}\sum_{\Lc'\in S^+_{i,e}(\Lc)} W_{\Lc'}(U)
\]
\[
\hspace{3cm}
+\sum_{i:e\in D_i(\Lc)}\sum_{\Lc'\in M^-_{U,i,e}(\Lc)} W_{\Lc'}(U)-\sum_{i:e\in D_i(\Lc)}\sum_{\Lc'\in M^+_{U,i,e}(\Lc)} W_{\Lc'}(U).
\]
Multiplying this identity by $\prod_{p\in P(\Lambda)}\chi_{\alpha_p}(U_{\partial p})$ and
integrating over $\U(N)^{E(\Lambda)}$ gives the expected formula.
\end{proof}

\end{document}
